\documentclass[11pt]{article}

\usepackage[a4paper,margin=1in]{geometry}
\usepackage[T1]{fontenc}
\usepackage{lmodern}
\usepackage{microtype}
\usepackage{amsmath,amssymb,amsthm,mathtools}
\usepackage{booktabs,tabularx}
\usepackage{enumitem}
\usepackage[numbers,sort&compress]{natbib}
\usepackage[colorlinks=true,linkcolor=blue,citecolor=blue,urlcolor=blue]{hyperref}
\usepackage[nameinlink,noabbrev]{cleveref}

\setlength{\emergencystretch}{3em}

\newtheorem{theorem}{Theorem}[section]
\newtheorem{proposition}[theorem]{Proposition}
\newtheorem{corollary}[theorem]{Corollary}
\newtheorem{lemma}[theorem]{Lemma}
\theoremstyle{definition}
\newtheorem{definition}[theorem]{Definition}
\theoremstyle{remark}
\newtheorem{remark}[theorem]{Remark}

\newcommand{\one}{\mathbf 1}
\newcommand{\Lzero}{\mathrm{L0}}
\newcommand{\Lone}{\mathrm{L1}}
\newcommand{\Ltwo}{\mathrm{L2}}

\title{\textbf{From Frozen Affine Pairs to Growing CRT Fibers}\\
\large Uniformity Envelopes, Cutoff Phase Diagrams, and the
Fixed-Two Quantifier Barrier}
\author{Liang Wang}
\date{July 2026}

\begin{document}
\maketitle

\begin{abstract}
TPC-137 proves full-M\"obius logarithmic cancellation for every
fixed determinant-two affine pair and fixed periodic coefficient.
The actual TPC family is not fixed.  This paper measures the gap by
comparing two squarefree decompositions.  The magnitude cutoff has
tail \(Y^{o(1)}(N/R+\sqrt Y)\), combined parameter modulus at most
\(MR_{\rm eff}^4\), and point census \(O(N+R_{\rm eff}^2)\).  The
prime-square cutoff has tail \(O(N/P+\sqrt Y)\), at most
\(4^{\pi(P)}\) ordered component pairs, and combined parameter
modulus
\[
 M\operatorname{lcm}(k,\ell)^2\le M(P^\#)^2.
\]
For every fixed integer data height \(Q\), Tao's fixed-affine
theorem is uniform over the corresponding finite record set.
A slow diagonal gives some non-effective
\(Q_{*,\omega}(x)\to\infty\), but no prescribed growth rate.
There is now also an explicit positive corridor:
Tao--Ter\"av\"ainen prove a power-of-log estimate, outside a small
logarithmic-density exceptional set, uniformly for positive affine
coefficients and constants of sufficiently small polylogarithmic
height.  This corrects the older ``fixed only'' frontier, but does
not contain an arbitrary actual CRT family and does not control
deterministic exceptional prefixes.  No ordinary all-prefix power
estimate or positive \(\Ltwo\) input is claimed.
\end{abstract}

\section{Two cutoff languages}

Let \(I\) be an integer interval of length \(N\), and let
\[
 D(z)=d+sz,\qquad V(z)=u+az,\qquad su-ad=2
\]
be primitive and positive on \(I\).  Suppose
\[
 \max_{z\in I}\{D(z),V(z)\}\le Y.
\]
Let \(M\) denote a fixed periodic modulus already present in the
coefficient.  The word ``combined modulus'' below refers to the
\(z\)-parameter.  Passing to the ambient \(n\)-frame may multiply
it by the separately recorded base modulus \(as\).

\subsection{Magnitude cutoff}

TPC-128 defines
\[
 Q_R(m)=\sum_{\substack{k\le R\\k^2\mid m}}\mu(k)
\]
and proves the two-mask tail
\begin{equation}
 \ll Y^{o(1)}\left(\frac NR+\sqrt Y\right),
\label{eq:magnitude-tail}
\end{equation}
as well as the point census \(O(N+R^2)\)
\citep{WangTPC128}.  Only
\[
 R_{\rm eff}=\min\{R,\sqrt Y\}
\]
can occur.  After also resolving the period,
\begin{equation}
 L_z\le MR_{\rm eff}^4.
\label{eq:magnitude-modulus}
\end{equation}
The exact-cutoff regime \(R\ge\sqrt Y\) has zero truncation tail,
not merely the bound \eqref{eq:magnitude-tail}.

\subsection{Prime-square cutoff}

TPC-137 instead uses
\[
 \mathsf S_P(m)
 =
 \prod_{p\le P}(1-\one_{p^2\mid m})
 =
 \sum_{k\mid P^\#}\mu(k)\one_{k^2\mid m}
\]
\citep{WangTPC137}.  Its Boolean nature yields
\begin{equation}
 \sum_{z\in I}
 \left|
 \mu^2(D(z))\mu^2(V(z))
 -\mathsf S_P(D(z))\mathsf S_P(V(z))
 \right|
 \ll\frac NP+\sqrt Y.
\label{eq:prime-tail}
\end{equation}

\begin{proposition}[Prime-cutoff component geometry]
\label{prop:prime-geometry}
Before compatibility removes empty components, the prime cutoff has
at most
\begin{equation}
 4^{\pi(P)}
\label{eq:pair-count}
\end{equation}
ordered pairs \((k,\ell)\).  Every nonempty pair, after resolving
the period \(M\), has parameter modulus
\begin{equation}
 L_z
 =
 \operatorname{lcm}(M,k^2,\ell^2)
 \le
 M\operatorname{lcm}(k,\ell)^2
 \le M(P^\#)^2.
\label{eq:prime-modulus}
\end{equation}
\end{proposition}

\begin{proof}
The primorial has \(2^{\pi(P)}\) divisors, independently in each
coordinate.  The identity
\(\operatorname{lcm}(k^2,\ell^2)
=\operatorname{lcm}(k,\ell)^2\) proves
\eqref{eq:prime-modulus}.
\end{proof}

\begin{remark}[The exponent correction]
The final upper bound in \eqref{eq:prime-modulus} contains
\((P^\#)^2\), not \((P^\#)^4\).  The latter would count the two
divisors as if their prime supports could be duplicated inside the
same squarefree primorial.
\end{remark}

\section{A fixed-height uniform envelope}

The fixed theorem of Tao cannot be called on an unspecified growing
record \citep{Tao2016LogChowla,WangTPC129}.  It does, however, give
uniformity over every finite bounded record set.

\begin{definition}[Finite affine record set]
\label{def:records}
For \(Q\ge1\), let \(\mathfrak D(Q)\) be the finite set of integer
records
\[
 \mathfrak d=(a_1,b_1,a_2,b_2,M,r)
\]
such that
\[
 1\le a_i,M\le Q,\quad |b_i|\le Q,\quad0\le r<M,
\]
the reduced forms are primitive and positive after a fixed
translation, and
\[
 0<|a_1b_2-a_2b_1|\le Q.
\]
The residue \(r\bmod M\) represents one periodic component.
\end{definition}

For an admissible terminal ratio \(\omega=\omega(x)\), define
\[
 \mathcal E_{\mathfrak d}(x,\omega)
 =
 \frac1{\log\omega}
 \left|
 \sum_{x/\omega<z\le x}
 \frac{
 \lambda(a_1z+b_1)\lambda(a_2z+b_2)
 \one_{z\equiv r\ ({\rm mod}\ M)}
 }z
 \right|.
\]

\begin{theorem}[Fixed-height uniformity]
\label{thm:fixed-height}
For every fixed \(Q\),
\begin{equation}
 \max_{\mathfrak d\in\mathfrak D(Q)}
 \mathcal E_{\mathfrak d}(x,\omega(x))
 \longrightarrow0
\label{eq:fixed-height}
\end{equation}
as \(x\to\infty\), whenever
\(1\le\omega(x)\le x\) and \(\omega(x)\to\infty\).
\end{theorem}

\begin{proof}
Resolve the one fixed residue in each record.  After factoring fixed
contents, the two forms remain nonparallel.  Tao's theorem gives
convergence for every individual record.  The maximum in
\eqref{eq:fixed-height} is over the finite set
\(\mathfrak D(Q)\).
\end{proof}

\begin{remark}[What finiteness does not give]
The convergence threshold in \cref{thm:fixed-height} may depend on
\(Q\).  Taking \(Q=Q(x)\) inside the theorem is invalid unless one
has first proved a uniform rate controlling that dependence.
\end{remark}

\section{A proved small-polylog almost-scale corridor}

The post-2025 source frontier is stronger than the fixed theorem.
The following is the affine special case recorded after
Tao--Ter\"av\"ainen's quantitative correlation theorem
\citep[Theorem~3.1 and Remarks~3.2]{TaoTeravainen2026}.

\begin{theorem}[Tao--Ter\"av\"ainen affine corridor]
\label{thm:tt-corridor}
There is a sufficiently small absolute \(c_0>0\) and an exceptional
set of scales \(\mathcal E\) such that
\begin{equation}
 \frac1{\log X}\int_{[1,X]\cap\mathcal E}\frac{dt}{t}
 \ll(\log X)^{-c_0},
\label{eq:tt-exceptional}
\end{equation}
and, for every \(N\notin\mathcal E\),
\begin{equation}
 \sup_{\substack{
 1\le a_i,b_i\le(\log N)^{c_0}\\
 a_1b_2\ne a_2b_1}}
 \frac1N\left|
 \sum_{n\le N}
 \lambda(a_1n+b_1)\lambda(a_2n+b_2)
 \right|
 \ll(\log N)^{-c_0}.
\label{eq:tt-affine}
\end{equation}
The constant \(c_0\) may be decreased so that one exponent controls
both the data height and the saving.
\end{theorem}

\begin{remark}[Source quantifier lock]
\label{rem:tt-source-lock}
The underlying \(\mathcal L\)-parameter theorem is local in a
separate ambient variable \(X\): it produces
\(\mathcal E_X\subset[\sqrt X,X]\), controls
\(N\in[\sqrt X,X]\setminus\mathcal E_X\), and is uniform only for
\[
 W\le\mathcal L^c,\qquad
 b,h_1,h_2=O(\mathcal L^c).
\]
Its non-pretentious branch assumes, in the notation of the source,
\[
 \exp M(g_1;X^2,\log^{1/125}X)\gg\mathcal L.
\]
The source verifies this condition for \(\lambda\) and then records
the global affine special case stated in
\cref{thm:tt-corridor}, using \(W=a_1a_2\) and shrinking the
absolute exponent.  We invoke that recorded special case, not an
extension to arbitrary periodic weights or squarefree masks.
\end{remark}

\begin{remark}[Exact eligibility test]
\label{rem:tt-eligibility}
Equation \eqref{eq:tt-affine} is callable only after every residue
split and content extraction has produced positive integer
coefficients and constants within its small-polylog height, with
nonzero reduced determinant.  A
dyadic block uses two prefix endpoints; its exceptional set must
therefore include the corresponding endpoint dilates.  Resolving
many CRT moduli also requires the union of their pulled-back
exceptional sets to retain vanishing logarithmic density.  None of
these checks follows merely from saying that the original
determinant is two.
\end{remark}

\begin{corollary}[Eligible finite-family transport]
\label{cor:tt-family}
Suppose a scale-\(N\) decomposition has total absolute coefficient
mass at most
\((\log N)^{\kappa_{\rm cen}+o(1)}\),
every reduced record satisfies \cref{rem:tt-eligibility}, and the
union of the required pulled-back exceptional sets has logarithmic
density \(o(1)\).  If \(\kappa_{\rm cen}<c_0\), then the retained
Liouville part has a power-of-log saving at every remaining scale.
\end{corollary}

\begin{proof}
Apply \eqref{eq:tt-affine} to every eligible component and sum
absolute values.  The resulting exponent is
\(c_0-\kappa_{\rm cen}+o(1)>0\).
\end{proof}

This corollary is a source-safe \(\Lone\) corridor, not an actual
TPC certificate: the literal data-height, exceptional-set union,
squarefree tail, weights, phases and prefixes still have to be
verified.

\section{A qualitative diagonal beyond the explicit corridor}

\begin{lemma}[Non-effective affine diagonal]
\label{lem:diagonal}
For each fixed admissible terminal-ratio function \(\omega\), there
exists a nondecreasing integer function
\(Q_{*,\omega}(x)\to\infty\) such that
\begin{equation}
 \max_{\mathfrak d\in\mathfrak D(Q_{*,\omega}(x))}
 \mathcal E_{\mathfrak d}(x,\omega(x))
 \longrightarrow0.
\label{eq:diagonal}
\end{equation}
No prescribed lower bound on \(Q_{*,\omega}(x)\) follows, and the
construction need not be uniform in \(\omega\).
\end{lemma}

\begin{proof}
For each integer \(j\), use \cref{thm:fixed-height} to choose
\(X_j\) so large that the maximum over \(\mathfrak D(j)\) is at
most \(1/j\) for \(x\ge X_j\).  Increase the \(X_j\) if necessary
and set \(Q_{*,\omega}(x)=j\) on \([X_j,X_{j+1})\).  The thresholds
come from qualitative convergence for the chosen \(\omega\), so
this construction emits neither an effective formula nor a named
growth class.
\end{proof}

\begin{corollary}[Exact containment interface]
\label{cor:containment}
Suppose every retained record in an actual family at scale \(x\)
has encoded height at most \(Q_{*,\omega}(x)\), and suppose its complete
absolute reassembly cost is \(x^{o(1)}\) with a rate compatible
with \eqref{eq:diagonal}.  Then the fixed-height conclusions can be
summed on that family.  Neither premise is supplied by
\cref{lem:diagonal}.
\end{corollary}

The compatibility clause is necessary: an unspecified \(x^{o(1)}\)
factor can grow faster than the unspecified decay in
\eqref{eq:diagonal}.

\section{Cutoff phase diagram}

The two decompositions serve different goals.

\begin{center}
\small
\begin{tabularx}{\textwidth}{@{}
>{\raggedright\arraybackslash}p{0.19\textwidth}
>{\raggedright\arraybackslash}p{0.23\textwidth}
>{\raggedright\arraybackslash}p{0.26\textwidth}
>{\raggedright\arraybackslash}X@{}}
\toprule
Cutoff & Elementary tail & Parameter modulus & Main obstruction\\
\midrule
Magnitude \(R\) &
\(Y^{o(1)}(N/R+\sqrt Y)\) &
\(MR_{\rm eff}^4\) &
divisor maximum, weighted census, growing affine data\\
Prime \(P\) &
\(N/P+\sqrt Y\) &
\(M(P^\#)^2\) &
\(4^{\pi(P)}\) components and growing data\\
\bottomrule
\end{tabularx}
\end{center}

For example, \(P\asymp\log\log x\) gives only a logarithmic tail.
More precisely, its \(1/P\) relative term then tends to zero only
like \(1/\log\log x\); the separate square-root term must also be
controlled.  On the positive side, the crude bounds
\[
 4^{\pi(P)}\le4^P,\qquad P^\#\le P^P
\]
show that its precompatibility costs are \(x^{o(1)}\).  Sharper
prime-number estimates show that for a sufficiently small constant
multiple of \(\log\log x\), the component count is
\((\log x)^{o(1)}\) and the primorial modulus is a small power of
\(\log x\).  This creates a genuine candidate for
\cref{thm:tt-corridor}, conditional on the literal base height and
all exceptional-set/reassembly checks.

Conversely, forcing the \(1/P\) relative part of the elementary
prime tail to be \(O(x^{-\delta})\) requires \(P\ge x^\delta\).
Then the cutoff and its record family lie far outside the
small-polylog corridor.  The effective value cap
\(k,\ell\le\sqrt Y\) can reduce the realized family, but it does
not by itself return it to that corridor.

\begin{definition}[Growing determinant-two affine gate]
\label{def:growing-gate}
A callable growing theorem must specify a function
\(\varepsilon_{\rm aff}(x;Q)\) and prove, in its own statement,
\[
 \sup_{\mathfrak d\in\mathfrak D(Q(x))}
 \mathcal E_{\mathfrak d}(x,\omega(x))
 \le\varepsilon_{\rm aff}(x;Q(x))
\]
for the declared \(Q(x)\), origins and windows.  An actual H3
return additionally needs the physical weights, phases and prefix
quantifiers, or a separately proved transfer for them.
\end{definition}

\section{Stops and claim boundary}

\paragraph{FIXED-DATA STOP.}
If any coefficient, residue modulus, period or cutoff grows outside
a proved envelope, the fixed theorem is not callable.

\paragraph{DIAGONAL STOP.}
The existence of \(Q_{*,\omega}(x)\) does not show that a prescribed actual
height is at most \(Q_{*,\omega}(x)\).  Renaming the slow diagonal
as the physical family is forbidden.  The diagonal is not the same
object as the explicit almost-scale corridor in
\cref{thm:tt-corridor}.

\paragraph{POWER STOP.}
The qualitative \(o(1)\) in \cref{thm:fixed-height} has zero
\(x\)-power exponent.  It cannot be entered as \(1/400\), nor can
it pay a positive physical loss.

\begin{center}
\small
\begin{tabularx}{\textwidth}{@{}p{0.49\textwidth}X@{}}
\toprule
Statement & Level and status\\
\midrule
Two cutoff identities, tails, moduli and component counts
& Exact/elementary \(\Lzero\).\\
Fixed-height uniform envelope
& Unconditional frozen \(\Lone\).\\
Small-polylog affine envelope outside a small exceptional scale set
& Unconditional sourced \(\Lone\); actual eligibility open.\\
Existence of some slow \(Q_{*,\omega}(x)\)
& Exact non-effective diagonal.\\
Larger prescribed polylogarithmic or power affine envelope
& Open beyond the sourced corridor.\\
Containment of the actual CRT family
& Not proved.\\
Ordinary all-prefix power saving or positive L2
& Open; no \(\Ltwo\) claim.\\
\bottomrule
\end{tabularx}
\end{center}

TPC-139 therefore locates the next analytic task precisely: prove
the growing gate of \cref{def:growing-gate}, or replace the
termwise CRT route.  Even an almost-scale theorem would still need
a selector theorem before it could control deterministic physical
prefixes, which is the subject of the next paper.

\begingroup
\footnotesize
\raggedright
\bibliographystyle{plainnat}
\bibliography{references}
\endgroup

\end{document}
